\section{AVALIAÇÃO DO PROJETO DO CURSO}

A avaliação do curso de \NOMECURSO deve ocorrer de forma contínua e dialogada através da interação de professores, gestores e estudantes. Ao longo do desenvolvimento das atividades curriculares, todos os segmentos, docentes, discentes e técnicos-administrativos, juntamente com a Coordenação do Curso, NDE e Colegiado do curso, devem atuar na direção da consolidação de mecanismos que possibilitem a permanente avaliação dos objetivos do curso.

\subsection{Avaliação do Docente}

O corpo docente que atua no curso passa por avaliações semestrais através de questionários respondidos pelos estudantes, por meio do sistema eletrônico Q-acadêmico, para cada componente curricular nos quais estão regularmente matriculados no semestre letivo. 

No instrumento são observados pontos como Pontualidade, Assiduidade, Domínio de conteúdo, Metodologia de Ensino, Avaliação e Relação Professor-Aluno.

A partir dos resultados obtidos são gerados relatórios para cada docente, na sua área de acesso no Q-Acadêmico, referentes a cada componente curricular, que devem ser percebidos pelo professor como instrumentos de autoavaliação e reflexão sobre a sua prática docente, avaliando-a e, se for o caso, intervir de forma a enriquecê-la e viabilizar melhorias na aprendizagem dos estudantes.


\subsection{Núcleo Docente Estruturante}


A constituição, funcionamento e atuação do NDE do \NOMECURSO estão em consonância com o disposto na Resolução CONSUP Nº 004, de 28 de janeiro de 2015, que determina a organização do Núcleo Docente Estruturante no IFCE, como sendo:
\begin{itemize}
	\item Construir e acompanhar a execução do PPC;
	\item Promover a revisão e atualização do PPC, tendo como principal objetivo a adequação do perfil profissional do egresso, devendo as alterações serem aprovadas pela maioria do NDE, e submetidas à análise e aprovação do colegiado do curso;
	\item Analisar os resultados obtidos nas avaliações internas e externas (ENADE, Relatório de Avaliação para Reconhecimento de curso) e propor estratégias para o desenvolvimento da qualidade acadêmica do curso;
	\item Zelar pelo cumprimento das Diretrizes Curriculares Nacionais para os cursos de graduação.
\end{itemize}

\subsection{Colegiado}


A constituição, funcionamento e atuação do colegiado do \NOMECURSO estão em consonância com o disposto na Resolução N° 75, de 13 de agosto de 2018, que define as normas de funcionamento dos colegiados de curso do IFCE.

Conforme estabelece o Artigo 4° da referida Resolução, compete ao Colegiado do IFCE:

\begin{itemize}
	\item Supervisionar as atividades curriculares, propondo aos órgãos competentes as medidas necessárias à melhoria do ensino, pesquisa e extensão;
	\item Aprovar as propostas de estruturação e reestruturação do Projeto Pedagógico do Curso;
	\item Avaliar o desenvolvimento do Projeto Pedagógico do Curso no tocante a sua atualização, primando pela sintonia com as demandas da sociedade e do mundo do trabalho;
	\item Deliberar sobre as recomendações propostas pelos docentes, discentes e egressos sobre assuntos de interesse do curso;
	\item Propor soluções para as questões administrativas e pedagógicas do curso, tais como aquelas que tratam de evasão, reprovação, retenção, entre outras;
	\item Propor, conforme o caso, a flexibilização curricular, bem como a extinção e a alteração de componentes curriculares;
	\item Coletar e analisar informações sobre as diferentes áreas do saber que compõem o curso, incluindo questões de cunho acadêmico;
	\item Orientar acerca de qual perfil docente deve ser solicitado, por ocasião de concurso público e/ou de remoção de professores, vislumbrando as necessidades do curso e as características de seu Projeto Pedagógico;
	\item Organizar e construir a sequência de afastamento docente no âmbito do curso, bem como deliberar acerca da efetivação deste afastamento, com base na regulamentação vigente;
	\item Colaborar, sempre que solicitado, no auxílio, indicação e escolha de membros de banca de concurso público, junto à Comissão Coordenadora de Concurso da Instituição;
	\item Receber, analisar e encaminhar demandas do corpo docente e discente e tomar decisões de natureza didático-pedagógicas sobre elas, desde que atendam à legislação em vigor.
\end{itemize}

\subsection{Encontros Pedagógicos}

Os Encontros Pedagógicos são realizados semestralmente, como uma ação formativa, em ambientes produtivos onde são preconizadas as práticas docentes e sua melhoria, cujo objetivo é provocar no professor a avaliação de sua prática docente para que ao longo do processo melhorem sua atuação pedagógica.

\subsection{Ações decorrentes dos processos de autoavaliação e avaliação externa}

De acordo com a Lei nº 10.861 de 14 de abril de 2004, que institui o Sistema Nacional de Avaliação da Educação Superior - SINAES, com o objetivo de assegurar processo nacional de avaliação das instituições de educação superior, a avaliação visa melhorar a qualidade da educação superior. Para isso, a avaliação institucional dos cursos e do desempenho dos alunos são mecanismos básicos para ponderações acerca da execução dos processos educativos na Educação Superior.

A Comissão Própria de Avaliação - CPA, do IFCE, instituída com base no art. 11 da lei nº 10.861/2004, tem a finalidade de implementar o processo de autoavaliação do Instituto, bem como a sistematização e a prestação das informações solicitadas pela Comissão Nacional de Avaliação da Educação Superior (CONAES). Esta comissão é formada por quatro representantes do corpo docente; quatro representantes do corpo técnico-administrativo; quatro representantes dos alunos; e quatro representantes da sociedade civil organizada, de acordo com o Regimento Interno da CPA do IFCE.

A autoavaliação institucional, tarefa da CPA, compõe-se da autoavaliação dos campi, que se dá em torno das seguintes informações: identificação e histórico do campus, identificação da subcomissão de avaliação e seus trabalhos e os resultados da autoavaliação por campus, com suas respectivas dimensões. 

Os instrumentos utilizados são questionários compostos de perguntas em cada dimensão descrita, aplicados com professores, alunos e técnicos. Os resultados são categorizados pelas respostas às perguntas em fragilidade, avaliação mediana ou potencialidade de cada grupo aplicado. Para conclusão, a classificação final se dá pela fragilidade, potencialidade, controvérsia ou tendência a um dos dois primeiros conceitos. 

A autoavaliação institucional prediz tomadas de decisões acerca do funcionamento do campus em torno de suas atividades, caracterizadas pelas dimensões avaliadas. Também representa a comunidade acadêmica e em geral na participação da melhoria do campus nos seus processos de criação e manutenção de mecanismos e estratégias corroboradoras de uma educação adequada à realidade local e global. 

Já a avaliação do curso, de caráter externo, compreende a avaliação do desempenho do estudante por meio do ENADE e das condições de ensino oferecidas a ele, especificamente o perfil do corpo docente, as instalações físicas e a organização didático-pedagógica. 

A periodicidade máxima de aplicação do ENADE é trienal, e consiste em contextualizar o perfil dos estudantes avaliados, o que fornece maiores subsídios para políticas e tomadas de decisões sobre o percurso do curso e da instituição na oferta do mesmo. O ENADE é componente curricular obrigatório nos cursos de graduação, sendo o mesmo explicitado aos estudantes pelo NDE, Colegiado e CPA. 

A avaliação externa, do SINAES, compreende o referencial básico para os processos de autorização, reconhecimento e renovação de reconhecimento de cursos de graduação. Portanto, o presente projeto vislumbra não somente a regulação do \NOMECURSO, como projeta a sua qualidade e estima diante da comunidade externa na participação de processos avaliativos que ensejem tal qualidade.

As avaliações internas e externas são importantes, pois a gestão do curso é realizada considerando a autoavaliação institucional e o resultado das avaliações externas como insumo para aprimoramento contínuo do planejamento do curso, com evidência da apropriação dos resultados pela comunidade acadêmica e existência de processo de autoavaliação periódica do curso.

Importante destacar que os dados coletados nas avaliações externas e da Comissão Própria de avaliação – CPA serão utilizadas como subsídio à Coordenadoria de Curso, à Coordenadoria Técnico-Pedagógica, à Diretoria de Ensino e aos próprios docentes para intervirem de forma a viabilizar melhorias no processo de ensino-aprendizagem. 

A partir do exame ENADE, por exemplo, é possível avaliar a qualidade do \NOMECURSO ofertado pelo IFCE Campus Maracanaú e o rendimento de seus discentes em relação aos conteúdos programáticos, a suas habilidades e competências. 

A partir das avaliações externas é possível verificar todos os tópicos contidos no Projeto Pedagógico do Curso e confrontar o que está escrito com as práticas docentes, infraestrutura, etc. 

Questões e resultados obtidos por meio da CPA, desempenho dos estudantes no ENADE, reconhecimento do curso e renovação do reconhecimento, serão pautados e discutidos amplamente pelo Colegiado do curso, divulgados à comunidade acadêmica e, sempre que necessário, serão feitos os devidos encaminhamentos ao NDE. 
